% 第六章 重点板块深度分析（三）：左侧布局型

\chapter{重点板块深度分析（三）：左侧布局型}

\section{医疗器械：政策边际改善+技术创新驱动，估值底部蓄力}

医疗器械板块自2021年中以来经历了近三年的深度调整，申万医疗器械指数最大回撤超过60\%，估值水平从峰值的60倍PE回落至当前的20倍左右，处于近十年15\%分位以下。我们认为，随着集采政策边际趋缓、医疗设备更新政策落地、国产高端设备突破加速，行业基本面已出现明确的边际改善信号，当前位置具备显著的左侧布局价值。

\subsection{景气度三维验证：业绩筑底+政策催化+国产突破三重共振}

\textbf{业绩端：收入增速逐季回升，盈利质量稳步改善。} 2025年申万医疗器械板块（剔除异常值）实现营业收入约4,200亿元，同比增长12.3\%；归母净利润约580亿元，同比增长8.7\%。进入2026年Q1，板块营收增速提升至15.2\%，归母净利润增速回升至14.6\%，呈现逐季改善态势。分细分领域看，医学影像设备增速最快（2026Q1同比+22\%），其次是体外诊断（+17\%）和心血管器械（+14\%），低值耗材受集采影响增速相对较慢（+5--8\%）。毛利率方面，板块整体毛利率维持在45--50\%区间，高端设备领域毛利率可达60\%以上，盈利质量显著优于医药行业整体水平。

\textbf{产业端：医疗新基建延续+设备更新周期启动。} 医疗新基建方面，2025年全国医疗卫生机构固定资产投资完成额约1.2万亿元，同比增长11.5\%，其中县级医院建设和区域医疗中心建设是主要抓手。医疗设备更新方面，2025年10月国家卫健委等部门联合印发《全面推进医疗设备更新实施方案》，明确提出2026--2027年完成全国范围内医疗机构设备更新升级，涉及设备品类包括影像设备、手术设备、重症设备、检验设备等十大类，预计撬动设备更新需求超2,000亿元。从落地节奏看，2026年Q1各地设备更新贷款已发放超800亿元，设备招标金额同比增长35\%以上，政策红利正加速转化为实际订单。

\textbf{政策端：集采规则优化+创新支持加码。} 集采政策出现明确的边际改善信号：一是竞价规则从"唯低价是取"转向"质量优先、价格合理"，中标价格降幅趋缓；二是集采品类从高值耗材向低值耗材、检验试剂延伸，但对创新器械给予更长的市场独占期；三是DRG/DIP支付方式改革深化，推动医疗机构从"耗占比"管控转向"价值医疗"，优质国产设备性价比优势更加凸显。创新支持方面，NMPA创新医疗器械特别审查通道持续提速，2025年共有45个创新医疗器械获批上市，同比增长36\%，其中国产高端影像设备、电生理器械、介入瓣膜等领域突破尤为显著。

\begin{exhibitbox}[图 6-1：医疗器械板块营收与净利润增速（2022Q1--2026Q1）]
\centering
\vspace{0.3cm}
\begin{tikzpicture}
\begin{axis}[
    width=0.85\linewidth, height=5.5cm,
    ylabel={同比增速（\%）},
    xlabel={季度},
    xtick={1,2,3,4,5,6,7,8,9,10,11,12,13,14,15,16,17},
    xticklabels={22Q1,22Q2,22Q3,22Q4,23Q1,23Q2,23Q3,23Q4,24Q1,24Q2,24Q3,24Q4,25Q1,25Q2,25Q3,25Q4,26Q1},
    x tick label style={rotate=45, anchor=east, font=\scriptsize},
    ymin=-10, ymax=30,
    ymajorgrids=true,
    grid style={dashed, gray!30},
    legend pos=north east,
    legend style={font=\footnotesize},
    symbolic x coords={1,2,3,4,5,6,7,8,9,10,11,12,13,14,15,16,17},
    enlarge x limits=0.03,
]
\addplot[navy, thick, mark=circle, mark size=2pt] coordinates {
    (1, 22) (2, 18) (3, 12) (4, 8)
    (5, 5) (6, 3) (7, 2) (8, -2)
    (9, -1) (10, 4) (11, 7) (12, 9)
    (13, 10) (14, 11) (15, 13) (16, 12)
    (17, 15)
};
\addlegendentry{营收增速}

\addplot[accentblue, thick, mark=square, mark size=2pt] coordinates {
    (1, 28) (2, 22) (3, 15) (4, 10)
    (5, 7) (6, 2) (7, -3) (8, -8)
    (9, -5) (10, 2) (11, 6) (12, 8)
    (13, 9) (14, 10) (15, 12) (16, 9)
    (17, 15)
};
\addlegendentry{净利润增速}
\end{axis}
\end{tikzpicture}
\vspace{0.2cm}
\sourcenote{Wind，本研究团队整理测算}
\end{exhibitbox}

\subsection{预期差来源：市场对集采过度悲观，低估创新器械成长属性}

我们认为，当前市场对医疗器械板块的认知存在三大预期差，这也是板块估值修复的核心驱动因素。

\textbf{预期差一：集采政策影响已进入边际改善通道，而非持续恶化。} 市场普遍将集采视为医疗器械行业的长期压制因素，预期降价幅度和覆盖范围将持续扩大。但我们观察到，集采政策已出现明确的优化调整：一是降价幅度趋缓，以心脏支架为例，第一轮集采平均降幅93\%，而后续批次耗材集采平均降幅已收窄至50--60\%，部分品类甚至仅30--40\%；二是中标规则优化，从"最低价中标"转向"多家中标、分层竞价"，头部企业凭借质量和服务优势更易中标；三是创新豁免机制逐步建立，对上市未满3年的创新器械暂不纳入集采。我们判断，集采对行业的边际冲击最大的阶段已经过去，后续政策环境将持续改善。

\textbf{预期差二：国产高端设备突破速度远超市场预期，进口替代空间巨大。} 市场普遍认为国产医疗器械仅在中低端领域有竞争力，高端设备仍依赖进口。但事实上，近年来国产高端设备突破速度显著加快：以CT、MR为代表的医学影像设备，国产品牌已掌握核心部件技术，3.0T超导MRI、64排以上CT等高端机型陆续获批上市，2025年国产CT在三级医院的渗透率已从2020年的不足10\%提升至25\%以上；超声设备方面，国产品牌在中高端市场的份额已突破30\%；内窥镜领域，国产硬镜和软镜均实现关键技术突破。我们测算，高端医疗器械整体国产化率仍不足20\%，未来5--10年进口替代空间广阔。

\textbf{预期差三：出海正成为国产器械企业的第二增长曲线。} 市场对医疗器械企业的增长预期主要基于国内市场，对海外市场贡献估计不足。实际上，头部器械企业已加快国际化布局：迈瑞医疗海外收入占比已达45\%，产品进入全球190多个国家和地区；联影医疗高端影像设备成功进入美国、欧洲、日本等发达市场，2025年海外收入同比增长超过60\%；南微医学内镜耗材在欧美市场份额持续提升。我们认为，中国医疗器械企业凭借"工程师红利+成本优势+快速迭代"，有望复制家电、手机等行业的出海路径，在全球市场占据重要地位。

\begin{keyinsight}[医疗器械板块投资逻辑：三重修复共振]
医疗器械板块当前正经历"估值修复+业绩修复+预期修复"三重共振。估值层面，板块PE从60倍回落至20倍，处于历史底部区域，安全边际充足；业绩层面，收入和利润增速已触底回升，2026年有望进入加速增长通道；预期层面，集采政策边际改善、国产高端设备突破、出海加速三大因素正持续修正市场悲观预期。我们判断，医疗器械板块已进入左侧布局的黄金窗口期。
\end{keyinsight}

\subsection{核心催化与时间窗口}

医疗器械板块的核心催化因素主要集中在以下三个方面，时间窗口覆盖2026年全年：

\textbf{催化一：医疗设备更新政策落地见效（2026年Q2--Q3）。} 医疗设备更新贷款发放和设备招标采购将在2026年二、三季度进入高峰期，相关企业订单和收入有望逐季超预期。重点关注影像设备、ICU设备、检验设备等更新需求大的领域。

\textbf{催化二：新一轮耗材集采结果落地（2026年Q3--Q4）。} 新一轮高值耗材集采预计在2026年下半年开标，若降价幅度符合或低于市场预期，将成为板块估值修复的重要催化剂。重点关注集采规则的优化方向和中标价格降幅。

\textbf{催化三：高端设备出海突破（持续兑现）。} 国产高端影像设备、手术机器人等在海外市场的注册和销售进展，是板块长期成长逻辑的重要验证点。若有重大海外订单或注册获批，将显著提升市场对行业长期成长空间的预期。

\subsection{产业链图谱与核心标的}

医疗器械产业链可划分为上游核心零部件、中游设备与耗材制造、下游医疗机构三个环节。上游高端零部件（如探测器、超导磁体、芯片等）技术壁垒最高，目前仍以进口为主，国产替代空间大；中游是行业核心，涵盖医学影像、体外诊断、高值耗材、低值耗材等多个细分领域；下游以公立医院为主体，集中度较低。

\begin{exhibitbox}[图 6-2：医疗器械产业链图谱与核心标的]
\centering
\vspace{0.4cm}
\begin{tikzpicture}[node distance=0.45cm, every node/.style={font=\small}, scale=1.0, transform shape]

% 上游
\node[draw, fill=navy!20, rectangle, rounded corners=2pt, minimum height=0.8cm, minimum width=2.8cm, align=center] (part1) at (0,0) {\textbf{核心零部件}\\ 探测器/磁体/芯片};
\node[draw, fill=navy!20, rectangle, rounded corners=2pt, minimum height=0.8cm, minimum width=2.8cm, align=center, right=of part1] (part2) {\textbf{原材料}\\ 医用塑料/金属/试剂原料};
\node[above=0.1cm of $(part1.north)!0.5!(part2.north)$, font=\bfseries\color{navy}] {上游：零部件与原材料};

% 中游
\node[draw, fill=accentblue!25, rectangle, rounded corners=2pt, minimum height=0.9cm, minimum width=2.6cm, align=center, below=1.0cm of part1] (mid1) {\textbf{医学影像设备}\\ 联影医疗/迈瑞/东软};
\node[draw, fill=accentblue!25, rectangle, rounded corners=2pt, minimum height=0.9cm, minimum width=2.6cm, align=center, right=of mid1] (mid2) {\textbf{体外诊断}\\ 迈瑞医疗/安图/新产业};
\node[draw, fill=accentblue!25, rectangle, rounded corners=2pt, minimum height=0.9cm, minimum width=2.6cm, align=center, right=of mid2] (mid3) {\textbf{高值耗材}\\ 惠泰医疗/南微/乐普};

\node[draw, fill=accentblue!20, rectangle, rounded corners=2pt, minimum height=0.8cm, minimum width=2.6cm, align=center, below=0.5cm of mid1] (mid4) {\textbf{生命信息支持}\\ 迈瑞医疗/理邦};
\node[draw, fill=accentblue!20, rectangle, rounded corners=2pt, minimum height=0.8cm, minimum width=2.6cm, align=center, right=of mid4] (mid5) {\textbf{低值耗材}\\ 奥美医疗/振德};
\node[draw, fill=accentblue!20, rectangle, rounded corners=2pt, minimum height=0.8cm, minimum width=2.6cm, align=center, right=of mid5] (mid6) {\textbf{家用医疗}\\ 鱼跃医疗/三诺};

% 中游标签（基于所有中游节点居中）
\node[fit=(mid1) (mid2) (mid3) (mid4) (mid5) (mid6), inner sep=0pt] (mid-group) {};
\node[above=0.1cm of mid-group, font=\bfseries\color{accentblue}] {中游：设备与耗材制造};

% 下游
\node[draw, fill=riskgreen!20, rectangle, rounded corners=2pt, minimum height=0.9cm, minimum width=3.2cm, align=center, below=1.0cm of mid4] (down1) {\textbf{公立医院体系}\\ 三级医院/二级医院/基层};
\node[draw, fill=riskgreen!20, rectangle, rounded corners=2pt, minimum height=0.9cm, minimum width=3.2cm, align=center, right=of down1] (down2) {\textbf{民营医疗机构}\\ 体检中心/专科/诊所};
\node[draw, fill=riskgreen!20, rectangle, rounded corners=2pt, minimum height=0.9cm, minimum width=2.6cm, align=center, right=of down2] (down3) {\textbf{零售与出口}\\ OTC零售/海外市场};
\node[above=0.1cm of down2, font=\bfseries\color{riskgreen}] {下游：终端市场};

% 箭头（mid1→mid4→down1 级联，避免箭头穿节点）
\draw[-{Stealth}, thick, steelgrey] (part1.south) -- (mid1.north);
\draw[-{Stealth}, thick, steelgrey] (part2.south) -- (mid2.north);
\draw[-{Stealth}, thick, steelgrey] (mid1.south) -- (mid4.north);
\draw[-{Stealth}, thick, steelgrey] (mid4.south) -- (down1.north);
\draw[-{Stealth}, thick, steelgrey] (mid3.south) -- (down2.north);
\draw[-{Stealth}, thick, steelgrey] (mid6.south) -- (down3.north);

\end{tikzpicture}
\vspace{0.3cm}
\sourcenote{本研究团队整理绘制}
\end{exhibitbox}

\textbf{核心标的分析：}

\begin{itemize}[leftmargin=2em]
\item \textbf{迈瑞医疗（300760）}：国内医疗器械绝对龙头，生命信息支持、体外诊断、医学影像三大业务协同发展。2025年营收约380亿元，净利润约110亿元。海外收入占比45\%，国际化布局领先。当前PE（2026E）约20倍，处于历史估值底部，具备显著的安全边际和长期成长空间。

\item \textbf{联影医疗（688271）}：国产高端医学影像设备龙头，CT、MR、PET-CT等全线产品布局，技术实力对标GPS（通用电气、飞利浦、西门子）。2025年营收突破120亿元，净利润约25亿元。高端产品加速进口替代，海外市场拓展迅猛，是国产高端医疗设备崛起的核心代表。

\item \textbf{开立医疗（300633）}：国产超声设备和内窥镜龙头，高端彩超和软式内窥镜技术国内领先。2025年营收约35亿元，净利润约6.5亿元。内窥镜业务进入爆发期，海外市场拓展顺利，有望成为继迈瑞、联影之后的又一器械出海标杆。

\item \textbf{南微医学（688029）}：内镜微创诊疗器械龙头，产品覆盖消化、呼吸、肿瘤等多个领域。2025年营收约25亿元，净利润约4.5亿元。海外收入占比超50\%，是国内少数成功进入欧美高端市场的器械企业之一。

\item \textbf{惠泰医疗（688617）}：电生理和血管介入器械龙头，国产电生理设备市占率第一。2025年营收约18亿元，净利润约5亿元。受益于集采中标份额提升和创新产品放量，业绩增长确定性高。
\end{itemize}

\begin{exhibitbox}[表 6-1：医疗器械核心标的盈利预测与估值对比]
\centering
\vspace{0.3cm}
\small
\begin{tabularx}{\textwidth}{@{}X r r r r r r@{}}
\toprule
\textbf{公司} & \textbf{\makecell[r]{市值\\（亿元）}} & \textbf{\makecell[r]{2025A营收\\（亿元）}} & \textbf{\makecell[r]{2025A\\增速}} & \textbf{\makecell[r]{2026E净利润\\（亿元）}} & \textbf{\makecell[r]{2026E\\PE}} & \textbf{\makecell[r]{PEG}} \\
\midrule
迈瑞医疗 & 2,450 & 380 & 15\% & 122 & 20.1 & 1.3 \\
联影医疗 & 1,580 & 122 & 25\% & 32 & 49.4 & 2.5 \\
开立医疗 & 280 & 35 & 22\% & 8.2 & 34.1 & 1.8 \\
南微医学 & 240 & 25 & 18\% & 5.6 & 42.9 & 2.4 \\
惠泰医疗 & 350 & 18 & 30\% & 6.8 & 51.5 & 1.8 \\
\cmidrule(lr){1-7}
板块平均 & --- & --- & 15\% & --- & 28.5 & 1.9 \\
\bottomrule
\end{tabularx}
\vspace{0.2cm}
\sourcenote{Wind一致预期，本研究团队整理，市值数据截至2026年6月20日}
\end{exhibitbox}

\subsection{风险因素与下行空间}

\textbf{集采降价超预期风险。} 若新一轮耗材集采降价幅度显著超过市场预期（如平均降幅超过70\%），将对相关企业的盈利能力造成较大冲击，可能导致板块估值进一步承压。但我们判断，随着集采规则优化，大幅超预期降价的概率较低。

\textbf{研发失败风险。} 高端医疗器械研发周期长、投入大、失败率高，若核心在研产品临床进展不及预期或注册失败，将对企业长期成长逻辑造成负面影响。

\textbf{医院采购放缓风险。} 若宏观经济下行导致地方财政压力加大，或医保控费力度超预期，可能影响医疗机构的设备采购意愿和采购预算，导致行业需求不及预期。

\textbf{海外贸易壁垒风险。} 若美国、欧洲等主要海外市场加强对中国医疗器械的准入限制或加征关税，将影响企业的国际化进程和海外收入增长。

\vspace{0.2cm}

从下行空间看，医疗器械板块当前估值已处于历史底部区域，且业绩增速已见底回升，我们判断大幅下跌的空间有限。以迈瑞医疗为例，20倍PE对应约7\%的收益预期，考虑15\%左右的业绩增速，下行空间约10--15\%，而长期上行空间可达50\%以上，风险收益比具备吸引力。

\vspace{0.5cm}

\section{动力电池龙头：价格战进入尾声，格局优化龙头受益}

动力电池板块自2022年见顶以来经历了两年多的调整，核心矛盾在于市场对产能过剩和价格战的过度担忧。我们认为，经过两年多的行业出清，价格战已进入尾声，行业供需格局正在边际改善，头部企业凭借技术、成本、客户和规模优势，市场份额持续提升，盈利能力触底回升，当前位置具备较高的左侧布局价值。

\subsection{景气度三维验证：需求高增+出清加速+盈利企稳}

\textbf{业绩端：收入维持高增长，盈利触底回升。} 2025年全球动力电池出货量约980GWh，同比增长38\%；中国动力电池出货量约680GWh，同比增长42\%，占全球比例约70\%。A股动力电池板块（宁德时代、比亚迪等头部企业）2025年合计营收约1.8万亿元，同比增长22\%；归母净利润约680亿元，同比增长3\%。净利润增速显著低于营收增速，主要原因是价格战导致产品价格大幅下降。但进入2026年Q1，板块营收增速维持在25\%左右，净利润增速回升至12\%，呈现"量增、价稳、利升"的积极态势，显示行业盈利底部已过。

\textbf{产业端：产能出清加速，格局持续优化。} 经过两年多的价格战，行业产能出清明显加速。一方面，二三线电池企业普遍陷入亏损，产能扩张大幅放缓，部分企业已退出市场；另一方面，头部企业凭借规模效应和技术优势，市场份额持续提升。2026年Q1，全球动力电池CR5市占率已从2023年的70\%提升至78\%，中国市场CR5市占率更是超过85\%。产能利用率方面，头部企业维持在85\%以上，而二三线企业普遍低于50\%，行业分化加剧，龙头优势不断扩大。

\textbf{政策端：海外政策有压有托，国内支持持续。} 海外方面，美国IRA法案、欧盟碳边境调节机制（CBAM）等贸易壁垒增加了中国电池企业出海的难度，但同时也加速了国内企业的海外产能布局。国内方面，新能源汽车购置税减免政策延续、动力电池回收体系建设加快、储能产业支持政策加码，为行业需求增长提供持续支撑。特别值得关注的是，储能电池需求爆发式增长，2025年全球储能电池出货量约320GWh，同比增长65\%，成为动力电池企业的第二增长曲线。

\begin{exhibitbox}[图 6-3：全球动力电池出货量与增速（2020--2026E）]
\centering
\vspace{0.3cm}
\begin{tikzpicture}
\begin{axis}[
    width=0.85\linewidth, height=5.5cm,
    ylabel={出货量（GWh）},
    xlabel={年份},
    xtick={2020,2021,2022,2023,2024,2025,2026},
    xticklabels={2020,2021,2022,2023,2024,2025,2026E},
    ymin=0, ymax=1400,
    ymajorgrids=true,
    grid style={dashed, gray!30},
    legend pos=north west,
    legend style={font=\footnotesize},
    symbolic x coords={2020,2021,2022,2023,2024,2025,2026},
    enlarge x limits=0.08,
]
\addplot[ybar, fill=navy!60, draw=navy, bar width=0.5cm] coordinates {
    (2020, 180) (2021, 320) (2022, 550) (2023, 700)
    (2024, 790) (2025, 980) (2026, 1250)
};
\addlegendentry{动力电池出货量}
\end{axis}
\begin{axis}[
    width=0.85\linewidth, height=5.5cm,
    ylabel={同比增速（\%）},
    axis y line*=right,
    axis x line=none,
    ymin=0, ymax=100,
    legend pos=north east,
    legend style={font=\footnotesize},
    symbolic x coords={2020,2021,2022,2023,2024,2025,2026},
    enlarge x limits=0.08,
]
\addplot[riskred, thick, mark=circle, mark size=2.5pt] coordinates {
    (2020, 25) (2021, 78) (2022, 72) (2023, 27)
    (2024, 13) (2025, 24) (2026, 28)
};
\addlegendentry{同比增速}
\end{axis}
\end{tikzpicture}
\vspace{0.2cm}
\sourcenote{SNE Research，GGII，本研究团队测算}
\end{exhibitbox}

\subsection{预期差来源：市场过度担忧产能过剩，忽视龙头份额提升}

我们认为，当前市场对动力电池板块的认知存在三大预期差，这是投资机会的核心来源。

\textbf{预期差一：市场担忧的"全面产能过剩"实际是"结构性过剩"，龙头企业产能并不过剩。} 市场普遍认为动力电池行业产能严重过剩，全行业将面临长期价格战。但我们认为，当前行业呈现典型的"结构性过剩"特征：低端产能严重过剩，二三线企业产能利用率不足50\%；而高端产能（如800V高压平台、大圆柱、半固态等技术路线）供给仍然偏紧，头部企业产能利用率维持在85\%以上。更重要的是，经过两年多的价格战，行业新增产能已大幅放缓，2025年行业新增产能同比下降40\%，预计2026年新增产能将进一步下降，供需格局有望在2026年底前后迎来实质性改善。

\textbf{预期差二：市场低估了储能等新应用对需求的拉动作用。} 市场对动力电池企业的增长预期主要基于新能源车需求，对储能等新应用的贡献估计不足。实际上，储能电池正在成为动力电池企业的重要增长极：宁德时代储能业务收入占比已达20\%，2025年同比增长超过80\%；比亚迪、亿纬锂能等企业储能业务增速也远超动力电池业务。我们测算，2026年全球储能电池需求将达到480GWh，同比增长50\%，成为电池行业需求增长的重要引擎。此外，户用储能、便携式储能、车网互动（V2G）等新兴应用场景也在快速拓展，持续打开行业成长空间。

\textbf{预期差三：市场对海外贸易壁垒的担忧过度，中国企业的全球竞争力仍在增强。} 市场普遍担心美国IRA法案、欧盟反补贴调查等贸易壁垒会严重影响中国电池企业的出海进程。但我们认为，中国电池企业的全球竞争力不仅来自成本优势，更来自技术领先、产业链完善和快速迭代能力。面对贸易壁垒，头部企业正在通过海外建厂、合资合作、技术授权等方式积极应对。宁德时代德国工厂已投产，美国福特合作项目进展顺利；比亚迪在泰国、巴西、匈牙利等地布局海外工厂；国轩高科、亿纬锂能等也在加速海外产能布局。我们判断，海外产能布局不会改变中国电池企业的全球竞争优势，反而会推动行业从"产品出口"向"产能+技术输出"升级。

\begin{keyinsight}[动力电池行业投资逻辑：从"贝塔行情"到"阿尔法行情"]
动力电池行业已从高速成长期的"贝塔行情"转向成熟期的"阿尔法行情"。过去行业普涨的格局难以再现，未来投资机会将更多集中在头部企业。核心逻辑在于：一是行业出清加速，龙头份额持续提升，"赢家通吃"特征日益显著；二是技术路线迭代加速，头部企业凭借研发优势不断拉大与二三线企业的差距；三是海外市场打开第二增长曲线，全球化能力强的企业将获得估值溢价。我们判断，当前正是左侧布局动力电池龙头的较好时机。
\end{keyinsight}

\subsection{核心催化与时间窗口}

动力电池板块的核心催化因素主要包括：

\textbf{催化一：碳酸锂价格企稳回升（2026年Q2--Q3）。} 碳酸锂价格自2022年高点的60万元/吨回落至2025年底的约8万元/吨，跌幅超过85\%。我们判断，随着上游锂矿产能出清和下游需求复苏，碳酸锂价格有望在2026年二、三季度企稳回升至10--12万元/吨区间。锂价企稳将有助于稳定电池企业的盈利预期，缓解市场对价格战的担忧。

\textbf{催化二：新技术路线量产落地（2026年Q3--Q4）。} M3P电池、半固态电池、大圆柱电池等新技术路线预计在2026年下半年开始批量装车，将成为板块重要的技术催化。新技术不仅带来产品性能提升和ASP提高，也将拉开头部企业与二三线企业的技术差距，强化龙头的竞争优势。

\textbf{催化三：海外建厂项目落地（持续兑现）。} 头部企业海外工厂投产和订单落地是行业长期成长逻辑的重要验证。重点关注宁德时代北美项目、比亚迪欧洲和东南亚工厂、国轩高科美国项目等的进展。若海外产能顺利投产并获得大客户订单，将显著提升市场对企业长期成长空间的预期。

\subsection{产业链图谱与核心标的}

动力电池产业链包括上游原材料、中游电池制造和下游应用三个环节。上游包括锂、钴、镍等有色金属和正极、负极、隔膜、电解液等四大主材；中游是动力电池电芯和PACK制造，是产业链的核心环节；下游包括新能源汽车、储能、两轮车等应用场景。

\begin{exhibitbox}[图 6-4：动力电池产业链图谱与核心标的]
\centering
\vspace{0.3cm}
\begin{tikzpicture}[node distance=0.4cm, every node/.style={font=\small}, scale=1.0, transform shape]

% 上游资源
\node[draw, fill=navy!20, rectangle, rounded corners=2pt, minimum height=0.7cm, minimum width=2.2cm, align=center] (res1) at (0,0) {\textbf{锂矿}\\ 赣锋锂业/天齐锂业};
\node[draw, fill=navy!20, rectangle, rounded corners=2pt, minimum height=0.7cm, minimum width=2.2cm, align=center, right=of res1] (res2) {\textbf{钴镍}\\ 华友钴业/格林美};
\node[draw, fill=navy!20, rectangle, rounded corners=2pt, minimum height=0.7cm, minimum width=2.2cm, align=center, right=of res2] (res3) {\textbf{石墨}\\ 贝特瑞/璞泰来};
\node[above=0.1cm of res2, font=\bfseries\color{navy}] {上游：金属资源};

% 上游材料
\node[draw, fill=accentblue!20, rectangle, rounded corners=2pt, minimum height=0.7cm, minimum width=2.2cm, align=center, below=0.8cm of res1] (mat1) {\textbf{正极材料}\\ 容百/当升/德方};
\node[draw, fill=accentblue!20, rectangle, rounded corners=2pt, minimum height=0.7cm, minimum width=2.2cm, align=center, right=of mat1] (mat2) {\textbf{负极材料}\\ 贝特瑞/璞泰来/杉杉};
\node[draw, fill=accentblue!20, rectangle, rounded corners=2pt, minimum height=0.7cm, minimum width=2.2cm, align=center, right=of mat2] (mat3) {\textbf{隔膜}\\ 恩捷股份/星源材质};
\node[draw, fill=accentblue!20, rectangle, rounded corners=2pt, minimum height=0.7cm, minimum width=2.2cm, align=center, below=0.4cm of mat2] (mat4) {\textbf{电解液}\\ 天赐材料/新宙邦};
\node[above=0.1cm of mat2, font=\bfseries\color{accentblue}] {上游：四大主材};

% 中游电池
\node[draw, fill=riskgreen!25, rectangle, rounded corners=2pt, minimum height=1.0cm, minimum width=3.0cm, align=center, below=1.0cm of mat4] (cell) {\textbf{动力电池制造（核心环节）}\\ 宁德时代 / 比亚迪 / 亿纬 / 国轩 / 欣旺达};
\node[above=0.1cm of cell, font=\bfseries\color{riskgreen}] {中游：电池制造};

% 下游
\node[draw, fill=riskamber!25, rectangle, rounded corners=2pt, minimum height=0.8cm, minimum width=2.5cm, align=center, below=0.9cm of cell] (app1) {\textbf{新能源汽车}\\ 特斯拉/比亚迪/新势力};
\node[draw, fill=riskamber!25, rectangle, rounded corners=2pt, minimum height=0.8cm, minimum width=2.5cm, align=center, right=of app1] (app2) {\textbf{储能系统}\\ 宁德时代/阳光/比亚迪};
\node[draw, fill=riskamber!25, rectangle, rounded corners=2pt, minimum height=0.8cm, minimum width=2.5cm, align=center, left=of app1] (app3) {\textbf{两轮车/其他}\\ 雅迪/小牛};
\node[above=0.1cm of app1, font=\bfseries\color{riskamber}] {下游：应用场景};

% 箭头（mat2→mat4→cell 级联，避免箭头穿节点）
\draw[-{Stealth}, thick, steelgrey] (res1.south) -- (mat1.north);
\draw[-{Stealth}, thick, steelgrey] (res2.south) -- (mat2.north);
\draw[-{Stealth}, thick, steelgrey] (res3.south) -- (mat2.north);
\draw[-{Stealth}, thick, steelgrey] (mat2.south) -- (mat4.north);
\draw[-{Stealth}, thick, steelgrey] (mat4.south) -- (cell.north);
\draw[-{Stealth}, thick, steelgrey] (cell.south) -- (app1.north);
\draw[-{Stealth}, thick, steelgrey] (cell.south) -- (app2.north);

\end{tikzpicture}
\vspace{0.2cm}
\sourcenote{本研究团队整理绘制}
\end{exhibitbox}

\textbf{核心标的分析：}

\begin{itemize}[leftmargin=2em]
\item \textbf{宁德时代（300750）}：全球动力电池绝对龙头，市占率超过35\%，技术、成本、客户优势全面领先。2025年营收约3,600亿元，净利润约520亿元。储能业务快速增长，海外布局领先。当前PE（2026E）约15倍，考虑其全球龙头地位和20\%+的长期增速，估值具备显著吸引力。

\item \textbf{比亚迪（002594）}：全球新能源汽车龙头，动力电池自产自用，垂直整合优势显著。刀片电池技术领先，外供比例持续提升。2025年汽车销量约450万辆，电池外供约80GWh。电池业务分拆后价值有望重估。

\item \textbf{亿纬锂能（300014）}：动力电池第二梯队龙头，消费电池+动力电池+储能电池多轮驱动。客户结构优质，海外大客户拓展顺利。2025年营收约750亿元，净利润约55亿元。储能业务增长迅速，大圆柱电池技术领先。

\item \textbf{国轩高科（002074）}：磷酸铁锂动力电池龙头，大众集团战略入股，海外客户拓展顺利。2025年营收约380亿元，净利润约18亿元。美国工厂项目进展顺利，有望成为北美市场重要玩家。

\item \textbf{欣旺达（300207）}：消费电池龙头，动力电池业务快速崛起。HEV电池市占率领先，动力电池客户覆盖多家主流车企。2025年营收约580亿元，净利润约22亿元。动力+储能业务进入收获期。
\end{itemize}

\begin{exhibitbox}[表 6-2：动力电池核心标的盈利预测与估值对比]
\centering
\vspace{0.3cm}
\small
\begin{tabularx}{\textwidth}{@{}X r r r r r r@{}}
\toprule
\textbf{公司} & \textbf{\makecell[r]{市值\\（亿元）}} & \textbf{\makecell[r]{2025A营收\\（亿元）}} & \textbf{\makecell[r]{2025A\\增速}} & \textbf{\makecell[r]{2026E净利润\\（亿元）}} & \textbf{\makecell[r]{2026E\\PE}} & \textbf{\makecell[r]{PEG}} \\
\midrule
宁德时代 & 7,850 & 3,620 & 22\% & 610 & 12.9 & 0.8 \\
比亚迪（整体） & 6,200 & 7,200 & 28\% & 480 & 12.9 & 0.6 \\
亿纬锂能 & 1,250 & 750 & 32\% & 72 & 17.4 & 0.9 \\
国轩高科 & 520 & 380 & 35\% & 28 & 18.6 & 0.9 \\
欣旺达 & 580 & 580 & 25\% & 32 & 18.1 & 1.1 \\
\cmidrule(lr){1-7}
板块平均 & --- & --- & 25\% & --- & 16.0 & 0.9 \\
\bottomrule
\end{tabularx}
\vspace{0.2cm}
\sourcenote{Wind一致预期，本研究团队整理，市值数据截至2026年6月20日}
\end{exhibitbox}

\subsection{风险因素与下行空间}

\textbf{价格战超预期风险。} 若行业产能出清进度慢于预期，或下游需求不及预期导致竞争加剧，可能引发新一轮价格战，影响行业整体盈利能力。但我们判断，头部企业盈利已触底，进一步大幅降价的空间和意愿有限。

\textbf{技术路线变革风险。} 动力电池技术路线迭代较快，若半固态、固态电池等新技术路线商业化进度超预期，可能对现有液态电池企业造成冲击。但头部企业在新技术方面布局较早，技术储备充足，受到的冲击相对较小。

\textbf{海外贸易壁垒风险。} 若美国、欧盟等主要市场进一步加强对中国电池产品的贸易限制，可能影响企业的海外收入增长和国际化进程。但头部企业已开始海外产能布局，可部分规避贸易壁垒风险。

\textbf{原材料价格大幅波动风险。} 锂、钴、镍等原材料价格大幅波动可能影响电池企业的成本控制和盈利稳定性。但经过前几年的波动，企业在原材料套期保值和长单锁价方面经验更加丰富，抗波动能力有所增强。

\vspace{0.2cm}

从下行空间看，动力电池龙头当前估值已处于历史底部区域，宁德时代PE仅13倍左右，显著低于其20\%+的长期业绩增速，PEG不到1，估值安全边际较高。我们判断，在行业基本面不出现重大恶化的前提下，板块下行空间约10--15\%，而长期上行空间可达50--100\%，具备较高的风险收益比。

\vspace{0.5cm}

\section{创新药龙头：国际化突破+医保谈判优化，困境反转可期}

创新药板块自2021年高点以来经历了近五年的深度调整，中证创新药指数最大回撤超过70\%，市场情绪极度低迷。我们认为，随着创新药出海突破加速、医保谈判规则优化、ADC等热门赛道持续兑现，行业基本面已出现边际改善信号，头部创新药企业正迎来困境反转的关键窗口期，具备显著的左侧布局价值。

\subsection{景气度三维验证：出海突破+支付改善+研发兑现}

\textbf{业绩端：收入稳步增长，研发投入持续加大。} 2025年A股创新药板块（重点覆盖公司）合计营业收入约2,800亿元，同比增长18\%；扣非净利润约320亿元，同比增长12\%。收入增长的主要驱动因素包括：创新药放量、海外授权收入增加、仿创结合企业的仿制药业务贡献稳定现金流。研发投入方面，板块研发费用率维持在15--20\%，头部企业研发费用率超过20\%，持续的高研发投入为长期成长奠定基础。值得关注的是，2025年国产创新药License-out交易金额达到历史新高，首付款+里程碑总额超过350亿美元，同比增长60\%以上，显示中国创新药的全球价值正得到越来越多的认可。

\textbf{产业端：创新管线厚度显著提升，热门赛道持续兑现。} 经过多年的研发积累，中国创新药企业的管线厚度和质量已大幅提升。从数量看，国内创新药临床申请数量连续多年保持30\%以上的增速，2025年IND申请数量超过1,200个，位居全球第二。从质量看，在ADC、双抗、基因治疗、细胞治疗等前沿领域，中国企业的研发水平已接近或达到国际先进水平。特别是ADC赛道，中国企业在研管线数量约占全球的40\%，多款产品展现出Best-in-Class潜力，成为中国创新药走向全球的先锋。

\textbf{政策端：医保谈判边际缓和，创新支持政策加码。} 医保谈判规则出现积极变化：一是新增药品平均降价幅度从过去的60\%左右收窄至2025年的约40\%，降价幅度明显趋缓；二是谈判续约规则更加透明可预期，简单续约不再大幅降价，给予企业更稳定的盈利预期；三是创新药准入速度加快，从获批到进入医保的时间从过去的2--3年缩短至1年左右，大大加速了创新药的放量。与此同时，国家持续加大对生物医药创新的支持力度，科创板、港股18A、北交所等多层次资本市场为创新药企业提供融资支持，CRO/CDMO产业生态日益完善，创新药研发的基础设施不断夯实。

\begin{exhibitbox}[图 6-5：中国创新药License-out交易金额与数量（2020--2025）]
\centering
\vspace{0.3cm}
\begin{tikzpicture}
\begin{axis}[
    width=0.85\linewidth, height=5.5cm,
    ylabel={交易金额（亿美元）},
    xlabel={年份},
    xtick={2020,2021,2022,2023,2024,2025},
    ymin=0, ymax=400,
    ymajorgrids=true,
    grid style={dashed, gray!30},
    legend pos=north west,
    legend style={font=\footnotesize},
    symbolic x coords={2020,2021,2022,2023,2024,2025},
    enlarge x limits=0.1,
]
\addplot[ybar, fill=navy!60, draw=navy, bar width=0.45cm] coordinates {
    (2020, 55) (2021, 130) (2022, 160) (2023, 185) (2024, 220) (2025, 350)
};
\addlegendentry{License-out交易金额（首付款+里程碑）}
\end{axis}
\begin{axis}[
    width=0.85\linewidth, height=5.5cm,
    ylabel={交易数量（项）},
    axis y line*=right,
    axis x line=none,
    ymin=0, ymax=80,
    legend pos=north east,
    legend style={font=\footnotesize},
    symbolic x coords={2020,2021,2022,2023,2024,2025},
    enlarge x limits=0.1,
]
\addplot[riskred, thick, mark=circle, mark size=2.5pt] coordinates {
    (2020, 18) (2021, 32) (2022, 42) (2023, 48) (2024, 58) (2025, 72)
};
\addlegendentry{交易数量}
\end{axis}
\end{tikzpicture}
\vspace{0.2cm}
\sourcenote{医药魔方，本研究团队整理}
\end{exhibitbox}

\subsection{预期差来源：市场低估创新药商业化能力与国际化价值}

我们认为，当前市场对创新药板块的认知存在三大预期差：

\textbf{预期差一：市场对创新药企业的商业化能力预期偏低，低估了创新药的放量速度。} 市场普遍认为国产创新药销售额普遍偏小，难以支撑高估值。但我们观察到，随着医保准入加速和企业商业化能力提升，国产创新药的放量速度正在显著加快。以PD-1抑制剂为例，恒瑞医药的卡瑞利珠单抗年销售额已突破80亿元，信达生物的信迪利单抗年销售额接近50亿元，远超市场初期预期。新一代创新药如ADC、双抗等，凭借更好的疗效和更高的价格，放量速度可能更快。我们测算，2025年销售额超过10亿元的国产创新药已有30多个，预计到2028年将增加到60个以上，"十亿级"大品种的梯队正在形成。

\textbf{预期差二：市场对创新药出海的价值认识不足，低估了国际化的成长空间。} 市场普遍将License-out视为"一锤子买卖"，对其长期价值估计不足。但我们认为，License-out只是中国创新药国际化的第一步，随着中国创新药质量的提升，未来将有越来越多的中国企业自主开展全球临床试验和商业化，实现从"授权出海"到"自主出海"的升级。以百济神州为例，其泽布替尼已在全球60多个国家获批上市，2025年全球销售额突破30亿美元，成为首个真正意义上的全球级国产创新药。我们判断，未来5--10年，中国创新药企业的海外收入占比将从当前的不足10\%提升至30\%以上，成长空间巨大。

\textbf{预期差三：市场对医保谈判的担忧过度，支付环境实际已出现边际改善。} 市场普遍将医保谈判视为创新药行业的长期压制因素，担心创新药进入医保后价格大幅下降、盈利空间有限。但我们认为，医保谈判的影响正在发生积极变化：一是降价幅度趋缓，2025年创新药医保谈判平均降幅约40\%，较前几年的60\%明显收窄；二是以价换量效果显著，创新药进入医保后销量通常有数倍增长，收入和利润反而可能增加；三是支付体系多元化，商保、惠民保、患者援助等支付方式快速发展，创新药的支付环境正在改善。我们判断，医保谈判对创新药企业盈利的边际冲击最大的阶段已经过去。

\begin{keyinsight}[创新药行业投资逻辑：从"泡沫化成长"到"高质量成长"]
中国创新药行业正在经历从"泡沫化成长"到"高质量成长"的转变。过去市场追捧靶点、管线数量等"虚"的指标，导致行业泡沫化严重；经过五年的调整，市场开始回归商业化能力、研发效率、国际化水平等"实"的指标，行业估值体系正在重构。我们认为，真正具备核心竞争力的头部创新药企业，已经度过了最困难的时期，正迎来困境反转的投资机会。
\end{keyinsight}

\subsection{核心催化与时间窗口}

创新药板块的核心催化因素主要包括：

\textbf{催化一：重磅品种海外获批（2026年Q2--Q4）。} 多款国产创新药正处于FDA/EMA注册的关键阶段，若成功获批将成为板块重要的催化剂。重点关注ADC、双抗等前沿领域的海外注册进展，以及百济神州、恒瑞医药、荣昌生物等企业的海外申报动态。

\textbf{催化二：新一轮医保谈判结果落地（2026年Q4）。} 2026年国家医保谈判预计在四季度进行，若降价幅度继续趋缓或续约规则进一步优化，将缓解市场对医保控费的担忧，成为板块估值修复的重要催化剂。

\textbf{催化三：ADC等热门赛道持续兑现（全年）。} ADC赛道是当前中国创新药最具全球竞争力的领域，相关产品的临床数据读出、海外授权、商业化进展等都可能成为板块的催化因素。重点关注DS-8201、TROP2 ADC、Claudin18.2 ADC等热门靶点的竞争格局和临床数据。

\subsection{产业链图谱与核心标的}

创新药产业链包括上游CRO/CDMO、中游创新药研发与生产、下游医药流通和医疗机构三个环节。上游CRO/CDMO为创新药研发提供外包服务，中国CRO/CDMO产业全球竞争力强；中游是产业链核心，包括Big Pharma和Biotech两类企业；下游包括医院、药店等销售终端。

\begin{exhibitbox}[图 6-6：创新药产业链图谱与核心标的]
\centering
\vspace{0.3cm}
\begin{tikzpicture}[node distance=0.45cm, every node/.style={font=\small}, scale=1.0, transform shape]

% 上游
\node[draw, fill=navy!20, rectangle, rounded corners=2pt, minimum height=0.7cm, minimum width=2.6cm, align=center] (up1) at (0,0) {\textbf{CRO}\\ 药明康德/康龙化成};
\node[draw, fill=navy!20, rectangle, rounded corners=2pt, minimum height=0.7cm, minimum width=2.6cm, align=center, right=of up1] (up2) {\textbf{CDMO}\\ 药明生物/凯莱英};
\node[draw, fill=navy!20, rectangle, rounded corners=2pt, minimum height=0.7cm, minimum width=2.6cm, align=center, right=of up2] (up3) {\textbf{实验动物/试剂}\\ 昭衍新药/泰坦科技};
\node[above=0.1cm of up2, font=\bfseries\color{navy}] {上游：研发生产外包};

% 中游
\node[draw, fill=accentblue!25, rectangle, rounded corners=2pt, minimum height=1.0cm, minimum width=3.0cm, align=center, below=1.0cm of up1] (mid1) {\textbf{Big Pharma}\\ 恒瑞医药/科伦药业\\ 中国生物制药};
\node[draw, fill=accentblue!25, rectangle, rounded corners=2pt, minimum height=1.0cm, minimum width=3.0cm, align=center, right=1.0cm of mid1] (mid2) {\textbf{Biotech}\\ 百济神州/信达生物\\ 荣昌生物/康方生物};
\node[above=0.1cm of mid1, xshift=2cm, font=\bfseries\color{accentblue}] {中游：创新药企业};

% 下游
\node[draw, fill=riskgreen!20, rectangle, rounded corners=2pt, minimum height=0.8cm, minimum width=2.8cm, align=center, below=1.0cm of mid1] (down1) {\textbf{医疗机构}\\ 三级医院/专科};
\node[draw, fill=riskgreen!20, rectangle, rounded corners=2pt, minimum height=0.8cm, minimum width=2.8cm, align=center, right=of down1] (down2) {\textbf{零售药店}\\ 益丰/大参林/国药};
\node[draw, fill=riskgreen!20, rectangle, rounded corners=2pt, minimum height=0.8cm, minimum width=2.8cm, align=center, right=of down2] (down3) {\textbf{支付方}\\ 医保/商保/惠民保};
\node[above=0.1cm of down2, font=\bfseries\color{riskgreen}] {下游：终端与支付};

% 箭头
\draw[-{Stealth}, thick, steelgrey] (up1.south) -- (mid1.north);
\draw[-{Stealth}, thick, steelgrey] (up2.south) -- (mid2.north);
\draw[-{Stealth}, thick, steelgrey] (mid1.south) -- (down1.north);
\draw[-{Stealth}, thick, steelgrey] (mid2.south) -- (down2.north);

\end{tikzpicture}
\vspace{0.2cm}
\sourcenote{本研究团队整理绘制}
\end{exhibitbox}

\textbf{核心标的分析：}

\begin{itemize}[leftmargin=2em]
\item \textbf{恒瑞医药（600276）}：国内创新药绝对龙头，研发投入和管线厚度行业第一。经历了仿制药集采和PD-1价格战的双重压力后，公司正进入新一轮创新药收获期，ADC、双抗、PROTAC等新一代产品管线丰厚，海外市场布局加速。2025年营收约260亿元，净利润约55亿元。当前PE（2026E）约30倍，处于历史估值底部，困境反转逻辑清晰。

\item \textbf{百济神州（688235）}：中国创新药国际化标杆，泽布替尼成为首个全球级国产创新药。公司研发投入大、管线质量高、商业化能力强，是国内最具全球竞争力的创新药企业。2025年全球产品收入约55亿美元，归母净亏损约5亿美元（亏损大幅收窄）。预计2026年有望实现盈亏平衡，进入利润释放期。

\item \textbf{信达生物（01801.HK）}：国内Biotech龙头企业之一，PD-1信迪利单抗销售额位居行业前列，生物类似药和创新药管线丰厚。公司从"follow"向"best-in-class/first-in-class"转型成效显著，IL-23、PCSK9等新产品贡献增量。2025年营收约55亿元，产品收入约50亿元。

\item \textbf{荣昌生物（688331）}：国产ADC龙头，维迪西妥单抗（RC48）是首个国产ADC药物，在HER2低表达胃癌、尿路上皮癌等适应症展现优异疗效。公司与Seagen达成全球独家授权，首付款2亿美元+里程碑24亿美元，是国产创新药License-out的标杆。新一代ADC管线丰厚，成长潜力大。

\item \textbf{科伦药业（002422）}：仿创结合的医药龙头，输液业务提供稳定现金流，创新药业务进入收获期。科伦博泰ADC管线全球领先，与默沙东达成多项总额超100亿美元的License-out合作。公司估值低廉，创新药业务价值尚未被充分认知。
\end{itemize}

\begin{exhibitbox}[表 6-3：创新药核心标的盈利预测与估值对比]
\centering
\vspace{0.3cm}
\small
\begin{tabularx}{\textwidth}{@{}X r r r r r r@{}}
\toprule
\textbf{公司} & \textbf{\makecell[r]{市值\\（亿元）}} & \textbf{\makecell[r]{2025A营收\\（亿元）}} & \textbf{\makecell[r]{2025A\\增速}} & \textbf{\makecell[r]{2026E净利润\\（亿元）}} & \textbf{\makecell[r]{2026E\\PE}} & \textbf{\makecell[r]{管线\\价值}} \\
\midrule
恒瑞医药 & 2,350 & 262 & 16\% & 78 & 30.1 & 高（ADC/双抗） \\
百济神州（A） & 2,100 & 400（约） & 45\% & -35（盈亏平衡） & N/A & 极高（全球管线） \\
信达生物（H） & 680 & 55 & 22\% & -5（减亏） & N/A & 高（自免/肿瘤） \\
荣昌生物（A） & 320 & 20 & 60\% & -2（减亏） & N/A & 高（ADC平台） \\
科伦药业 & 650 & 230 & 12\% & 28 & 23.2 & 高（ADC管线） \\
\cmidrule(lr){1-7}
板块平均 & --- & --- & 18\% & --- & 28.0 & --- \\
\bottomrule
\end{tabularx}
\vspace{0.2cm}
\sourcenote{Wind一致预期，本研究团队整理，市值数据截至2026年6月20日}
\end{exhibitbox}

\subsection{风险因素与下行空间}

\textbf{临床失败风险。} 创新药研发成功率低，核心管线临床失败是创新药企业面临的最大风险。若重磅产品临床数据不及预期或失败，可能导致公司估值大幅下挫。建议重点关注研发管线多样性高、抗风险能力强的头部企业。

\textbf{医保降价超预期风险。} 若医保谈判降价幅度显著超过市场预期，将对创新药企业的盈利能力造成较大冲击。但我们判断，随着医保谈判规则优化，大幅超预期降价的概率较低。

\textbf{FDA审批风险。} 国产创新药海外注册面临FDA审批不确定性，如临床数据质量、人种差异、生产质量等方面的问题可能导致审批延迟或被拒。建议重点关注已在全球多中心开展临床试验、国际化经验丰富的企业。

\textbf{研发投入回报不及预期风险。} 创新药企业研发投入巨大，若研发效率不高或管线商业化价值有限，可能导致投入产出比低下，影响长期回报。

\vspace{0.2cm}

从下行空间看，创新药板块经过近五年的深度调整，估值已处于历史底部区域，市场预期极低，利空因素已基本反映在股价中。以恒瑞医药为例，30倍PE对应创新药业务而言已处于合理偏低水平，考虑其管线价值和长期成长空间，下行空间约15--20\%。而一旦行业情绪回暖或基本面超预期改善，上行空间可达50--100\%。我们判断，当前位置左侧布局创新药龙头具备较高的赔率。

\vspace{0.5cm}

\needspace{6\baselineskip}
\begin{keyinsight}[左侧布局型板块投资策略总结]
医疗器械、动力电池龙头、创新药龙头三大板块共同特征是：基本面已出现边际改善信号，但市场预期仍然偏低，估值处于历史底部区域，具备较高的左侧布局价值。三个板块的左侧属性略有不同：医疗器械的改善确定性最高、弹性中等，适合稳健型投资者左侧布局；动力电池龙头的估值最便宜、盈利拐点最明确，适合追求确定性的价值型投资者；创新药龙头的弹性最大、不确定性也最高，适合风险承受能力较强的成长型投资者。

配置建议上，我们建议将三大左侧布局型板块作为组合的"卫星配置"，占比约15--20\%，与明星型板块的"核心配置"（占比40--50\%）形成互补。左侧布局的核心是"分批建仓、耐心持有"，不必追求买在最低点，而是在估值底部区域逐步建立仓位，等待基本面和情绪面的共振修复。
\end{keyinsight}

